\documentclass[10pt,a4paper]{amsart}
\usepackage[margin=19mm]{geometry}
\usepackage{amsmath,amssymb,mathtools}
\usepackage[scr=rsfs,cal=euler]{mathalfa}
\usepackage{xfrac}
\usepackage[unicode,hidelinks,bookmarks=false]{hyperref}
\newcommand{\ie}{i.e.\ }
\newcommand{\cf}{cf.\ }
\newcommand{\resp}{respectively\ }
\newcommand{\Subsection}{\subsection}
\providecommand{\nobreakdash}{\nobreak\hbox{-}}
\setlength{\emergencystretch}{2em}
\multlinegap0pt
\usepackage{mathtools}
\usepackage{enumitem}
\newtheorem{theorem}{Theorem}
\newtheorem{lemma}[theorem]{Lemma}
\newtheorem{proposition}[theorem]{Proposition}
\newtheorem{corollary}[theorem]{Corollary}
\newtheorem{remark}[theorem]{Remark}
\newtheorem{question}[theorem]{Question}
\DeclareMathOperator{\Br}{Br}
\DeclareMathOperator{\Pic}{Pic}
\DeclareMathOperator{\NS}{NS}
\DeclareMathOperator{\End}{End}
\DeclareMathOperator{\im}{im}
\newcommand{\OX}{\mathcal O_X}
\newcommand{\DD}{\mathcal D}
\newcommand{\Gm}{\mathbb G_m}
\begin{document}
\title[A note on Azumaya algebras and one-forms]{A note on Azumaya algebras and one-forms}
\author{Siqing Zhang}
\date{}
\maketitle
\noindent\emph{Source and licence.} A note on Azumaya algebras and one-forms. Version arXiv:2605.30279v1 (CC BY 4.0). This material is distributed under the Creative Commons Attribution 4.0 International licence, http://creativecommons.org/licenses/by/4.0/, as recorded in the arXiv OAI licence metadata for this exact version and shown on the abstract page https://arxiv.org/abs/2605.30279v1. Full rights notice: licenses/arxiv-2605.30279-CC-BY-4.0.txt.\par\smallskip
\noindent\textit{Editorial scope.} Counted unit: the original \texttt{theorem} environment \texttt{-} at source lines 187--190 together with its complete proof at lines 192--194; the delivered document keeps source lines 58--195 so that every referenced label and every citation resolves inside it. Native editable source is the paper's own concatenated TeX; the AMS wrapper is editorial formatting only. No publisher class, upstream script or unrelated artwork is redistributed.
\begin{question}\cite[Question 7.3(1)]{Petrov}\label{petrov-question}
Does there exist a smooth proper variety \(X\) together with a smooth
subvariety \(Z\subset T^*X'\), finite over \(X'\), such that
\(\mathcal A_X|_Z\) does not split?
\end{question}

In this note, we answer this question affirmatively in every positive characteristic.

\section{The Answer}

 Let
\(
X' := X\times_{k,F_k} k
\)
be the Frobenius twist.
Write \(F_X=F_{X/k}\colon X\to X'\) for the relative
Frobenius, and \(\pi_X\colon X'\to X\) for the projection.
Milne's exact sequence, in the form used by Ogus--Vologodsky
\cite[(4.1.1)]{OV}, is the following exact sequence of \'{e}tale sheaves of
abelian groups on \(X'\):
\begin{equation}\label{eq:milne}
0\longrightarrow \mathcal O_{X'}^*
\longrightarrow F_{X,*}\mathcal O_X^*
\xrightarrow{d\log}
F_{X,*}Z^1_{X/k}
\xrightarrow{\pi_X^*-C_X}
\Omega^1_{X'/k}
\longrightarrow 0.
\end{equation}
Here \(Z^1_{X/k}\subset \Omega^1_{X/k}\) is the sheaf of closed one-forms and
\(C_X\) is the Cartier operator.  Splice the exact sequence and the connecting morphisms define the obstruction homomorphism
\begin{equation}\label{eq:obstruction}
\operatorname{ob}_{X'}\colon
H^0(X',\Omega^1_{X'/k})
\to
H^2_{\mathrm{\acute et}}(X',\mathcal O_{X'}^*).
\end{equation}



Let $\DD_{X/k}$ be the sheaf of crystalline differential operators on $X$.
 The center
of \(F_{X,*}\DD_{X/k}\) is identified with
\(\operatorname{Sym}_{\mathcal O_{X'}}T_{X'/k}\), and 
\(F_{X,*}\DD_{X/k}\) is identified with an Azumaya algebra on \(T^*X'\).  This is the algebra
denoted \(\mathcal A_X\) above.

\begin{lemma}\label{lem:graph}
Let \(X/k\) be smooth, and let
\(\omega\in H^0(X',\Omega^1_{X'/k})\).  Let
\(
i_\omega\colon X'\longrightarrow T^*X'
\)
be the section corresponding to \(\omega\), with graph
\(\Gamma_\omega\subset T^*X'\).  Then
\[
[i_\omega^*\mathcal A_X]
=\operatorname{ob}_{X'}(\omega)
\quad\text{in } H^2_{\mathrm{\acute et}}(X',\mathcal O_{X'}^*).
\]
In particular, if \(\operatorname{ob}_{X'}(\omega)\ne 0\), then
\(\mathcal A_X|_{\Gamma_\omega}\) does not split.
\end{lemma}

\begin{proof}
Ogus--Vologodsky identify \(\operatorname{ob}_{X'}(\omega)\) with the class of
the \(\Gm\)-gerbe of line bundles on \(X\) with integrable connection and
\(p\)-curvature \(\omega\) \cite[Proposition 4.2]{OV}.  The same gerbe is the gerbe of splittings of the restriction of
\(\mathcal A_X\) to the section \(i_\omega\) by \cite[Remark 4.3]{OV}.  
\end{proof}

The following proposition reduces Question \ref{petrov-question} to finding a variety with non-closed global 1-forms.

\begin{proposition}\label{prop:dimension}
Let \(X/k\) be a smooth proper connected variety.  If not every global
one-form on \(X\) is closed, i.e.
\(
H^0(X,Z^1_{X/k})\subsetneq H^0(X,\Omega^1_{X/k}),
\)
then the obstruction map
\(
\operatorname{ob}_{X'}\) as in \eqref{eq:obstruction}
is nonzero.
\end{proposition}

\begin{proof}
Let \(G:=\Pic^{\nabla,\mathrm{rig}}_{X/k}\)
be the rigidified Picard scheme of line bundles with integrable
connection.  Let \(V:=H^0(X',\Omega^1_{X'/k})\), viewed as a vector group.
By \cite[Proposition 4.11 and the paragraph following it]{OV}, taking \(p\)-curvature defines a morphism of group schemes
\[
\psi\colon G\longrightarrow V.
\]
By \cite[Proposition 4.2]{OV}, \(\operatorname{ob}_{X'}(\omega)=0\) if and only if \(\omega\in\psi(G(k))\).  Thus it is enough to show that
\(\psi(G(k))\ne V(k)\).

By Cartier descent,
\[
\Pic_{X'/k}\longrightarrow G,\qquad
M\longmapsto (F_X^*M,\nabla_{\mathrm{can}})
\]
identifies \(\Pic_{X'/k}\) with \(\ker(\psi)\); see \cite[Theorem 5.1]{Katz}.
Forgetting the connection fits into the exact sequence of fppf sheaves
\begin{equation*}
0\to H^0(X,Z^1_{X/k})\to G\xrightarrow{b}\Pic_{X/k}
\to H^1(X,Z^1_{X/k}),
\end{equation*}
as in \cite[Proposition 4.11]{OV}.  In particular, if \(G^0\) is the neutral
component of \(G\), then
\[
\dim G^0\leq h^0(X,Z^1_{X/k})+\dim \Pic^0_{X/k}.
\]
Let \(\operatorname{Im}(\psi|_{G^0})\) be the algebraic subgroup image.
The quotient morphism \(G^0\to \operatorname{Im}(\psi|_{G^0})\) is 
faithfully flat. Therefore
\[
\dim \operatorname{Im}(\psi|_{G^0})=\dim G^0-\dim\ker(\psi|_{G^0})
\leq h^0(X,Z^1_{X/k})<\dim V.
\]


Finally, passing from \(G^0\) to \(G\) only gives finitely many translates:
Let \(G_1=b^{-1}(\Pic^0_{X/k})\).  Since \(G_1\) is of finite type,
\(G_1/G^0\) is finite.  Moreover \(G(k)/G_1(k)\) injects into \(\NS(X)\), so
its image under \(\psi\) is finite: \(\NS(X)\) is finitely generated, whereas
\(V(k)\) is killed by \(p\).

\end{proof}


\begin{theorem}
    Let $X/k$ be a smooth proper connected variety.
    If not every global one-form on $X$ is closed, then there is a smooth subvariety $Z\subset T^*X'$, finite over $X'$, such that $\mathcal{A}_X|_Z$ does not split. Moreover, the finite morphism $Z\to X'$ is an isomorphism.
\end{theorem}

\begin{proof}
    Combine Lemma \ref{lem:graph} and Proposition \ref{prop:dimension}, and take $Z$ to be $\Gamma_{\omega}$ for some $\omega$ such that $\operatorname{ob}_{X'}(\omega)\neq 0$.
\end{proof}


\begin{thebibliography}{99}
\bibitem[Katz]{Katz} N. M. Katz, \emph{Nilpotent connections and the monodromy theorem: applications of a result of Turrittin}, Publ. Math. IHES \textbf{39} (1970), 175--232.
\bibitem[OV]{OV} A. Ogus and V. Vologodsky, \emph{Nonabelian Hodge theory in characteristic \(p\)}, Publ. Math. IHES \textbf{106} (2007), 1--138.
\bibitem[Petrov]{Petrov} A. Petrov, \emph{Non-abelian Hodge theory in positive characteristic, after Ogus--Vologodsky}, notes for the workshop ``Complex and \(p\)-Adic Simpson Correspondence'', University of Maryland, November 2025, version of May 4, 2026. Available at \url{https://math.mit.edu/~apetrov/papers/pdf/Maryland_OV.pdf}.
\end{thebibliography}
\end{document}
